\documentclass[11pt]{article}

\usepackage[a4paper,margin=2.3cm]{geometry}
\usepackage{amsmath,amssymb,bm}
\usepackage{booktabs}
\usepackage{longtable}
\usepackage{array}
\usepackage[hidelinks]{hyperref}
\usepackage{xcolor}

\emergencystretch=3em
\sloppy

\newcommand{\pred}[1]{\widehat{#1}}
\newcommand{\code}[1]{\texttt{\small #1}}
\newcommand{\Real}{\mathbb{R}}
\newcommand{\relu}{\operatorname{relu}}
\newcommand{\sig}{\sigma}
\newcommand{\logp}{\operatorname{log1p}}
\DeclareMathOperator*{\amax}{max}

\title{Loss Functions of the Neural AC-OPF Surrogates\\
\large GridFM (\code{gridfm-graphkit}), GridSFM, and LUMINA\\[2pt]
\normalsize A code-grounded formulation of objectives, constraint handling, and cost treatment}
\author{}
\date{\today}

\begin{document}
\maketitle

\begin{abstract}
\noindent
This note writes out, as explicit formulas, the training objective of each of the three
neural AC-OPF surrogates available here --- \code{gridfm-graphkit}, \code{GridSFM}, and the
LUMINA SDK --- together with the exact mechanism by which each equality constraint (KCL),
each inequality constraint (voltage, generator, thermal, angle) and the generation cost
enters the pipeline. Every formula is traced to the source file and line range that
implements it, and where a mechanism is \emph{not} present in the shipped code that is
stated explicitly rather than inferred from the corresponding paper. Two structural facts
organise the comparison. First (Section~\ref{sec:common}): none of the three solves a KKT
system at inference; they are learned maps whose constraint satisfaction is either
\emph{architectural} (guaranteed by an output transformation) or \emph{penalised}
(encouraged by a training term, hence violable). Second, the three sit at very different
points on that axis --- GridSFM carries an explicit term for every constraint family and
for the cost, GridFM carries power-balance and a reactive-limit penalty but no cost term,
and the released LUMINA objective (Section~\ref{sec:lumina}) is a plain supervised
regression loss with no physics, no cost, and no Lagrangian at all.
\end{abstract}

\tableofcontents

% ======================================================================
\section{Notation and the classical reference problem}
\label{sec:notation}

Let $\mathcal{N}$ be the bus set, $\mathcal{G}$ the generator set, $\mathcal{E}$ the branch
set (AC lines and transformers), and $\mathcal{L}$ the load set. The classical AC-OPF is
\begin{align}
\min_{P_g,Q_g,V,\theta}\quad
& C(P_g) \;=\; \sum_{g\in\mathcal{G}}
\left( c_{2,g} P_{g}^{2} + c_{1,g} P_{g} + c_{0,g}\right)
\label{eq:opf-obj}\\
\text{s.t.}\quad
& P_{i}^{\mathrm{inj}}(V,\theta) \;=\; \textstyle\sum_{g\in\mathcal{G}_i} P_g - \sum_{l\in\mathcal{L}_i} P_{d,l} - G_{s,i}V_i^2,
&& \forall i\in\mathcal{N}
\label{eq:kcl-p}\\
& Q_{i}^{\mathrm{inj}}(V,\theta) \;=\; \textstyle\sum_{g\in\mathcal{G}_i} Q_g - \sum_{l\in\mathcal{L}_i} Q_{d,l} + B_{s,i}V_i^2,
&& \forall i\in\mathcal{N}
\label{eq:kcl-q}\\
& V_i^{\min}\le V_i\le V_i^{\max}, && \forall i\in\mathcal{N}
\label{eq:vlim}\\
& P_g^{\min}\le P_g\le P_g^{\max},\qquad Q_g^{\min}\le Q_g\le Q_g^{\max}, && \forall g\in\mathcal{G}
\label{eq:glim}\\
& |S_{ij}|\le S_{ij}^{\max},\qquad |S_{ji}|\le S_{ij}^{\max}, && \forall (i,j)\in\mathcal{E}
\label{eq:therm}\\
& \Delta\theta_{ij}^{\min}\le \theta_i-\theta_j \le \Delta\theta_{ij}^{\max}, && \forall (i,j)\in\mathcal{E}.
\label{eq:ang}
\end{align}
An interior-point solver (IPOPT via PowerModels.jl, or \code{pandapower}) introduces
multipliers $\lambda$ for \eqref{eq:kcl-p}--\eqref{eq:kcl-q} and $\mu\ge 0$ for
\eqref{eq:vlim}--\eqref{eq:ang} and drives the KKT residual of the Lagrangian
\begin{equation}
\mathcal{L}(x,\lambda,\mu)=C(P_g)+\lambda^\top h(x)+\mu^\top g(x)
\label{eq:kkt}
\end{equation}
to zero by Newton iteration, per instance. A solution of \eqref{eq:kkt} is denoted with a
star: $x^\star=(P_g^\star,Q_g^\star,V^\star,\theta^\star)$, with optimal cost
$C^\star=C(P_g^\star)$.

Predictions of a neural model carry a hat: $\pred{x}=(\pred{P}_g,\pred{Q}_g,\pred{V},\pred\theta)$.
For a batch of $G$ grids, $b(\cdot)$ maps an element (bus, generator, branch) to its graph
index, and for a per-element quantity $u$ we write the per-graph mean and per-graph max as
\begin{equation}
\langle u\rangle_k \;=\; \frac{1}{|\{e: b(e)=k\}|}\sum_{b(e)=k} u_e,
\qquad
[u]^{\amax}_k \;=\; \max_{b(e)=k} u_e ,
\qquad k=1,\dots,G .
\label{eq:pergraph}
\end{equation}
$\logp(u)=\log(1+u)$, and $\sig(\cdot)$ is the logistic sigmoid.

% ======================================================================
\section{The structure common to all three models}
\label{sec:common}

All three models learn a parametric map
\begin{equation}
f_\phi:\;\underbrace{(\text{topology},\,P_d,Q_d,\,\text{limits},\,\text{cost coefficients})}_{\text{graph input}}
\;\longmapsto\;
(\pred{P}_g,\pred{Q}_g,\pred{V},\pred\theta),
\label{eq:map}
\end{equation}
fitted offline against a dataset of solver solutions $\{x^{\star(n)}\}_{n=1}^{N}$. At
inference one forward pass replaces the Newton loop; no $\lambda,\mu$ are computed and no
KKT residual is driven to tolerance. Consequently constraint satisfaction arises only
through the four mechanisms
\begin{enumerate}
\item \textbf{imitation}: a supervised term pulling $\pred{x}$ towards $x^\star$, which is
      feasible, so $\pred{x}$ inherits approximate feasibility;
\item \textbf{penalty}: a differentiable residual term added to the loss --- \emph{soft},
      violable at inference;
\item \textbf{output transformation}: an algebraic map making a box constraint hold by
      construction --- \emph{hard}, up to the exact form used;
\item \textbf{post-hoc evaluation}: a violation metric computed after prediction, with no
      gradient path.
\end{enumerate}
The baseline supervised term shared in spirit by all three is
\begin{equation}
L_{\mathrm{sup}}
=\big\|\pred{P}_g-P_g^\star\big\|^2
+\big\|\pred{Q}_g-Q_g^\star\big\|^2
+\big\|\pred{V}-V^\star\big\|^2
+\big\|\pred\theta-\theta^\star\big\|^2 ,
\label{eq:sup}
\end{equation}
though each implementation reweights, normalises, masks and caps the four blocks
differently; the exact forms follow.

% ======================================================================
\section{GridFM (\code{gridfm-graphkit})}
\label{sec:gridfm}

\subsection{Task formulation: masked reconstruction}

GridFM poses PF, OPF and SE as \emph{masked reconstruction} on a heterogeneous graph with
node types \code{bus} and \code{gen}. A boolean mask $M$ marks the entries the model must
reconstruct; unmasked entries are clamped to ground truth inside the network
(\code{gnn\_heterogeneous\_gns.py:246--247}):
\begin{equation}
\pred{u} \;=\; M\odot u_{\mathrm{net}} \;+\; (1-M)\odot u^{\mathrm{given}} .
\label{eq:gridfm-clamp}
\end{equation}
The bus feature vector supplies $P_d,Q_d,Q_g,V_m,V_a$, the bus-type one-hots
$(\mathrm{PQ},\mathrm{PV},\mathrm{REF})$, the limits $V^{\min},V^{\max},Q_g^{\min},Q_g^{\max}$,
and the shunts $G_s,B_s$; the generator features supply $P_g,P_g^{\min},P_g^{\max}$ and the
cost coefficients $c_0,c_1,c_2$ (\code{datasets/globals.py:1--40}). The task therefore
changes only which columns are masked, not the architecture: for \code{PowerFlow} the loss
regresses $(V_m,V_a)$; for \code{OptimalPowerFlow} it regresses $(V_m,V_a,Q_g)$ at buses
plus $P_g$ at generators (\code{training/loss.py:143--148}).

\subsection{Architectural (hard) constraint handling}

Two box constraints are enforced by construction, and only in the OPF task
(\code{gnn\_heterogeneous\_gns.py:249--260}, using \code{models/utils.py:193--195}):
\begin{align}
\pred{V}_i &= V_i^{\min} + \big(V_i^{\max}-V_i^{\min}\big)\,\sig\!\left(z_{V,i}\right),
\label{eq:gridfm-vbound}\\
\pred{P}_g &= P_g^{\min} + \big(P_g^{\max}-P_g^{\min}\big)\,\sig\!\left(z_{P,g}\right),
\label{eq:gridfm-pbound}
\end{align}
so \eqref{eq:vlim} and the $P_g$ half of \eqref{eq:glim} hold strictly (open interval) for
every OPF prediction. In the PF task these two lines are skipped and no box constraint is
architectural.

Reactive generation is not a free output at all. The OPF physics decoder
(\code{models/utils.py:82--117}) \emph{derives} $\pred{Q}_g$ from the nodal reactive
balance at PV and REF buses,
\begin{equation}
\pred{Q}_{g,i}
=\underbrace{Q_{i}^{\mathrm{in}}(\pred{V},\pred\theta)}_{\text{aggregated branch flows}}
+\;Q_{d,i}\;-\;q^{\mathrm{shunt}}_i,
\qquad
q^{\mathrm{shunt}}_i = B_{s,i}\pred{V}_i^{2},
\qquad i\in \mathrm{PV}\cup\mathrm{REF},
\label{eq:gridfm-qg}
\end{equation}
and sets $\pred{Q}_{g,i}=0$ elsewhere. Equation~\eqref{eq:kcl-q} is thus satisfied
\emph{by construction} at PV/REF buses --- but this consumes the $Q_g$ degree of freedom,
so the $Q_g$ box \eqref{eq:glim} is \emph{not} guaranteed and needs the penalty of
Section~\ref{sec:gridfm-qviol}. The branch flows entering \eqref{eq:gridfm-qg} are the
exact $\pi$-model sending-end flows built from the stored $Y_{ff},Y_{ft}$ columns
(\code{models/utils.py:25--54}):
\begin{equation}
\pred{S}_{ij}=\pred{V}_i\,\overline{\left(Y_{ff,ij}\pred{V}_i+Y_{ft,ij}\pred{V}_j\right)},
\qquad
\pred{V}_i=\pred{V}_{m,i}e^{\mathrm{j}\pred{V}_{a,i}},
\qquad
P^{\mathrm{in}}_i=\!\!\sum_{(i,j)\in\mathcal{E}}\!\!\Re\{\pred{S}_{ij}\}.
\label{eq:gridfm-flows}
\end{equation}

\subsection{The loss terms in the registry}

The trainer builds a weighted sum over a configured list of registered losses
(\code{MixedLoss}, \code{training/loss.py:270--339}):
\begin{equation}
\boxed{\;
L_{\mathrm{GridFM}}=\sum_{m} w_m\, L_m ,
\;}
\label{eq:gridfm-mixed}
\end{equation}
with $\{(L_m,w_m)\}$ read from the YAML \code{training.losses} / \code{training.loss\_weights}.
The individual terms are as follows.

\paragraph{(a) Masked supervised terms.}
With $\mathcal{M}_{\mathrm{bus}}$ the masked bus entries over the task columns
$\mathcal{C}$ ($\mathcal{C}=\{V_m,V_a,Q_g\}$ for OPF, $\{V_m,V_a\}$ for PF) and
$\mathcal{M}_{\mathrm{gen}}$ the masked generator $P_g$ entries
(\code{training/loss.py:95--160}):
\begin{equation}
L_{\mathrm{bus}}=\frac{1}{|\mathcal{M}_{\mathrm{bus}}|}\!\!\sum_{(i,c)\in\mathcal{M}_{\mathrm{bus}}}\!\!\left(\pred{u}_{ic}-u^\star_{ic}\right)^2,
\qquad
L_{\mathrm{gen}}=\frac{1}{|\mathcal{M}_{\mathrm{gen}}|}\!\!\sum_{g\in\mathcal{M}_{\mathrm{gen}}}\!\!\left(\pred{P}_g-P^\star_g\right)^2 .
\label{eq:gridfm-masked}
\end{equation}
The unified variant \code{MaskedReconstructionMSE} (\code{training/loss.py:163--245})
replaces the pair by a single masked MSE over the six bus columns
$(V_m,V_a,P_g^{\mathrm{agg}},Q_g,P_d,Q_d)$, with the generator targets scattered onto buses,
$P_{g,i}^{\mathrm{agg}}=\sum_{g\in\mathcal{G}_i}P_g^\star$, and the bus-level $P_g$ mask set
where \emph{any} generator at that bus is masked.

\paragraph{(b) Layered physics residual --- the KCL penalty actually used.}
The network re-evaluates the power balance after \emph{every} message-passing layer
$\ell=0,\dots,L-1$. The per-bus residuals are (\code{models/utils.py:176--190})
\begin{align}
r^{P,(\ell)}_i &= \pred{P}^{(\ell)}_{g,i}-P_{d,i}+p^{\mathrm{shunt},(\ell)}_i-P^{\mathrm{in},(\ell)}_i,
& p^{\mathrm{shunt}}_i&=-G_{s,i}\pred{V}_i^2,
\label{eq:gridfm-rp}\\
r^{Q,(\ell)}_i &= \pred{Q}^{(\ell)}_{g,i}-Q_{d,i}+q^{\mathrm{shunt},(\ell)}_i-Q^{\mathrm{in},(\ell)}_i,
& q^{\mathrm{shunt}}_i&=+B_{s,i}\pred{V}_i^2,
\label{eq:gridfm-rq}
\end{align}
and are reduced to one scalar per layer (\code{gnn\_heterogeneous\_gns.py:295--299}):
\begin{equation}
R^{(\ell)}=\Big\langle \big\|(r^{P,(\ell)}_i,\,r^{Q,(\ell)}_i)\big\|_2 \Big\rangle_{i\in\mathcal{N}} .
\label{eq:gridfm-Rl}
\end{equation}
\code{LayeredWeightedPhysics} (\code{training/loss.py:342--381}) combines them with
normalised geometric weights of ratio $\beta=$ \code{base\_weight} (0.5 in the shipped
configs), so later layers dominate:
\begin{equation}
\boxed{\;
L_{\mathrm{phys}}=\sum_{\ell=0}^{L-1} \tilde w_\ell\, R^{(\ell)},
\qquad
\tilde w_\ell=\frac{\beta^{\,L-1-\ell}}{\sum_{m=0}^{L-1}\beta^{\,L-1-m}} .
\;}
\label{eq:gridfm-layered}
\end{equation}
The same residual is also fed back into the latent state,
$h^{(\ell)}_{\mathrm{bus}}\!\leftarrow\! h^{(\ell)}_{\mathrm{bus}}+\mathrm{MLP}(r^{P,(\ell)},r^{Q,(\ell)})$
(\code{gnn\_heterogeneous\_gns.py:300}), which is what makes this a physics-\emph{informed}
architecture rather than only a physics-penalised one.

\paragraph{(c) Full-Ybus power-balance loss (\code{PBE}, alternative).}
\code{PBELoss} (\code{training/loss.py:445--636}) assembles the complete sparse admittance
matrix from the edge attributes plus the shunt diagonal
$Y_{ii}\mathrel{+}= G_{s,i}+\mathrm{j}B_{s,i}$, then penalises the complex mismatch in
$L_1$:
\begin{equation}
\boxed{\;
L_{\mathrm{PBE}}
=\Big\langle\;\big| \pred{S}^{\mathrm{net}}_i - \pred{V}_i\,(\overline{Y}\,\overline{\pred{V}})_i \big|\;\Big\rangle_{i\in\mathcal{N}},
\qquad
\pred{S}^{\mathrm{net}}_i=(\pred{P}_{g,i}-\pred{P}_{d,i})+\mathrm{j}(\pred{Q}_{g,i}-\pred{Q}_{d,i}),
\;}
\label{eq:gridfm-pbe}
\end{equation}
each entry of $\pred{S}^{\mathrm{net}}$ again clamped to ground truth where unmasked, per
\eqref{eq:gridfm-clamp}. The real and imaginary parts are logged separately as the active
and reactive power loss in p.u.

\paragraph{(d) Reactive-limit penalty.}
\label{sec:gridfm-qviol}
\code{QgViolationPenalty} (\code{training/loss.py:639--685}) is a one-sided hinge on the
$Q_g$ box, averaged \emph{only over the violating buses}:
\begin{equation}
\boxed{\;
L_{Q_g\text{-viol}}
=\underbrace{\Big\langle \relu\!\big(\pred{Q}_{g,i}-Q^{\max}_{g,i}\big)\Big\rangle_{i:\,\pred{Q}_{g,i}>Q^{\max}_{g,i}}}_{\text{overshoot}}
+\underbrace{\Big\langle \relu\!\big(Q^{\min}_{g,i}-\pred{Q}_{g,i}\big)\Big\rangle_{i:\,\pred{Q}_{g,i}<Q^{\min}_{g,i}}}_{\text{undershoot}} ,
\;}
\label{eq:gridfm-qviol}
\end{equation}
with empty-set NaNs replaced by $0$. Note this is a \emph{mean over violators}, not a mean
over all buses: the gradient magnitude does not shrink as the number of violating buses
falls.

\paragraph{Instantiation used in the shipped OPF configs.}
For \code{HGNS\_OPF\_datakit\_case118.yaml} the objective is exactly
\begin{equation}
L_{\mathrm{GridFM}}^{\mathrm{OPF}}
= 0.1\,L_{\mathrm{phys}}
+ 0.1\,L_{\mathrm{gen}}
+ 0.75\,L_{\mathrm{bus}}
+ 0.001\,L_{Q_g\text{-viol}} ,
\label{eq:gridfm-instance}
\end{equation}
whereas the PF config uses $0.1\,L_{\mathrm{phys}}+0.9\,L_{\mathrm{bus}}$.

\subsection{Generation cost in GridFM}

There is \textbf{no cost term in any registered GridFM loss}. The cost coefficients
$c_0,c_1,c_2$ are input features, and the quadratic cost is evaluated only in
\code{tasks/opf\_task.py:84--102}, inside \code{test\_step}, as a reported metric:
\begin{equation}
\pred{C}_k=\!\!\sum_{g:\,b(g)=k}\!\!\left(c_{0,g}+c_{1,g}\pred{P}_g+c_{2,g}\pred{P}_g^2\right),
\qquad
\mathrm{gap}=\Big\langle \big|\pred{C}_k-C^\star_k\big|/C^\star_k\Big\rangle_k\times 100\% .
\label{eq:gridfm-gap}
\end{equation}
Cost therefore reaches the parameters \emph{only} through the labels $P^\star_g$ in
\eqref{eq:gridfm-masked}, which were produced by a cost-minimising solver. The precise
statement is: GridFM's forward pass runs no optimiser over $C$, and its training gradient
contains no $\partial C/\partial\pred{P}_g$ contribution; it learns the cost-optimal
dispatch by imitation. Thermal limits \eqref{eq:therm} and angle limits \eqref{eq:ang} are
likewise \emph{evaluation-only} in this codebase --- $\relu(|\pred{S}_{ft}|-S^{\max})$ is
computed in \code{test\_step} (\code{opf\_task.py:125--135}) and never in a training loss.

% ======================================================================
\section{GridSFM}
\label{sec:gridsfm}

GridSFM is the only one of the three whose released training objective contains an explicit
term for \emph{every} category: imitation, cost, KCL, branch flow, thermal loading, thermal
limit, physics stress, and feasibility. All of it lives in
\code{GridSFM/model/gridsfm/loss.py}.

\subsection{Predictions and architectural (hard) constraint handling}

The heads emit $\pred{x}=(\pred\theta,\pred{V})$ per bus, $(\pred{P}_g,\pred{Q}_g)$ per
generator, and one feasibility logit per graph. Box constraints are imposed by a
\emph{leaky clip} (\code{blocks.py:43--46}), applied with a mid-range offset for the
generator heads (\code{model.py:234--261}):
\begin{align}
\mathrm{lclip}(u;\ell,h)&=\operatorname{clamp}(u,\ell,h)+\alpha\Big[\relu(u-h)-\relu(\ell-u)\Big],
\label{eq:sfm-lclip}\\
\pred{P}_g&=\mathrm{lclip}\!\left(z_{P,g}+\tfrac12(P^{\min}_g+P^{\max}_g);\,P^{\min}_g,P^{\max}_g\right),
\label{eq:sfm-pg}\\
\pred{Q}_g&=\mathrm{lclip}\!\left(z_{Q,g}+\tfrac12(Q^{\min}_g+Q^{\max}_g);\,Q^{\min}_g,Q^{\max}_g\right),
\label{eq:sfm-qg}\\
\pred{V}_i&=\mathrm{lclip}\!\left(z_{V,i};\,V^{\min}_i,V^{\max}_i\right).
\label{eq:sfm-v}
\end{align}
With the default leak $\alpha=0$ (\code{model.py:63}) equations
\eqref{eq:sfm-pg}--\eqref{eq:sfm-v} enforce \eqref{eq:vlim} and \eqref{eq:glim}
\emph{exactly}; for $\alpha>0$ they are enforced up to an $\alpha$-scaled excursion that
keeps a gradient alive outside the box. The angle head is a DC-prior residual, anchored at
the slack bus (\code{model.py:243--248}):
\begin{equation}
\pred\theta = \operatorname{anchor}\!\Big(
\pi\tanh\!\big(\theta^{\mathrm{DC}}/\pi\big)
+ s\tanh\!\big(z_\theta/s\big)\Big),
\qquad s=\code{theta\_res\_scale},
\label{eq:sfm-theta}
\end{equation}
which bounds $\pred\theta$ globally but does \emph{not} enforce the per-branch angle
difference limits \eqref{eq:ang}.

\subsection{Branch flows and residuals}

Flows use the full $\pi$-model with off-nominal tap $\tau$ and phase shift $\phi$, in both
directions (\code{loss.py:185--209}); with $g+\mathrm{j}b=(r+\mathrm{j}x)^{-1}$ and
$\delta=\pred\theta_i-\pred\theta_j-\phi$:
\begin{align}
\pred{P}_{ij}&=\frac{g\,\pred{V}_i^2}{\tau^2}-\frac{\pred{V}_i\pred{V}_j}{\tau}\big(g\cos\delta+b\sin\delta\big),
&
\pred{Q}_{ij}&=-\frac{(b+b_{\mathrm{fr}})\pred{V}_i^2}{\tau^2}+\frac{\pred{V}_i\pred{V}_j}{\tau}\big(b\cos\delta-g\sin\delta\big),
\label{eq:sfm-flow-ij}\\
\pred{P}_{ji}&=g\,\pred{V}_j^2-\frac{\pred{V}_i\pred{V}_j}{\tau}\big(g\cos\delta-b\sin\delta\big),
&
\pred{Q}_{ji}&=-(b+b_{\mathrm{to}})\pred{V}_j^2+\frac{\pred{V}_i\pred{V}_j}{\tau}\big(b\cos\delta+g\sin\delta\big).
\label{eq:sfm-flow-ji}
\end{align}
The KCL residuals recompute the flows from $(\pred\theta,\pred{V})$ rather than reading the
separate edge head, deliberately, so that the term is a genuine physics check on the bus
state (\code{loss.py:265--317}):
\begin{align}
r^P_i &= \underbrace{\textstyle\sum_{j}\pred{P}_{ij}+\sum_j \pred{P}_{ji}}_{\text{flow out of }i}
\;-\;\Big(\pred{P}_{g,i}-P_{d,i}-G_{s,i}\pred{V}_i^2\Big),
\label{eq:sfm-rp}\\
r^Q_i &= \textstyle\sum_{j}\pred{Q}_{ij}+\sum_j \pred{Q}_{ji}
\;-\;\Big(\pred{Q}_{g,i}-Q_{d,i}+B_{s,i}\pred{V}_i^2\Big).
\label{eq:sfm-rq}
\end{align}

\subsection{The nine loss terms}

Write $\mathcal{F}=\{k:\ \text{feasible}_k=1\}$ for the feasible subset of the batch; all
regression terms are masked to $\mathcal{F}$ (\code{loss.py:97--119}), because the synthetic
wrapper injects infeasible samples whose ``labels'' are not solutions. Let
$\kappa=\code{PER\_ELEM\_CAP}=100$ and define the soft cap
$\mathrm{cap}(u)=\kappa\tanh(u/\kappa)$.

\paragraph{(1--4) Capped, normalised imitation.}
(\code{loss.py:690--728}) The angle error is taken on the circle and against the identically
anchored label; the voltage error is standardised by the per-graph label spread
$s_k=\max(\operatorname{std}_k(V^\star),0.01)$:
\begin{align}
L_\theta&=\Big\langle\Big\langle \mathrm{cap}\big(\operatorname{atan2}(\sin\Delta_i,\cos\Delta_i)^2\big)\Big\rangle_{i\in k}\Big\rangle_{k\in\mathcal{F}},
&&\Delta_i=\pred\theta_i-\operatorname{anchor}(\theta^\star)_i,
\label{eq:sfm-Ltheta}\\
L_V&=\Big\langle\Big\langle \mathrm{cap}\big((\pred{V}_i-V^\star_i)^2/s_{b(i)}^2\big)\Big\rangle_{i\in k}\Big\rangle_{k\in\mathcal{F}},
\label{eq:sfm-LV}\\
L_{P_g}&=\Big\langle\Big\langle \mathrm{cap}\big((\pred{P}_g-P^\star_g)^2\big)\Big\rangle_{g\in k}\Big\rangle_{k\in\mathcal{F}},
\label{eq:sfm-Lpq}\\
L_{Q_g}&=\Big\langle\Big\langle \mathrm{cap}\big((\pred{Q}_g-Q^\star_g)^2\big)\Big\rangle_{g\in k}\Big\rangle_{k\in\mathcal{F}} .
\label{eq:sfm-Lqg}
\end{align}
The $\tanh$ cap is what keeps a single pathological bus from owning the whole gradient.

\paragraph{(5) KCL, capacity-normalised with a worst-bus tilt.}
(\code{loss.py:590--631}) Residuals are normalised by the per-graph \emph{generation
capacity} $D^P_k=\sum_{g\in k}|P^{\max}_g|$, $D^Q_k=\sum_{g\in k}|Q^{\max}_g|$ --- not by
demand, so that the term is invariant to the load-spike perturbation modes used to
synthesise infeasible samples. With $\rho^P_i=|r^P_i|/D^P_{b(i)}$, margin
$\epsilon=0.01$ and tilt $w=0.8$:
\begin{equation}
\boxed{\;
L_{\mathrm{KCL},P}
=\Big\langle\;
w\,\big[\relu(\rho^P_i-\epsilon)\big]^{\amax}_k
+(1-w)\,\big\langle \rho^P_i\big\rangle_k
\;\Big\rangle_{k\in\mathcal{F}} ,
\;}
\label{eq:sfm-kcl}
\end{equation}
and identically for $L_{\mathrm{KCL},Q}$ with $\rho^Q$. The $80\%$ weight on the worst bus
means the term behaves closer to an $\ell_\infty$ than an $\ell_2$ penalty.

\paragraph{(6) Supervised branch flow.}
(\code{loss.py:531--552}) Against the solver's four-column edge label, in raw p.u.$^2$,
macro-averaged over edge types and graphs:
\begin{equation}
L_{\mathrm{br},P}=\Big\langle\big(\pred{P}_{ij}-P^\star_{ij}\big)^2+\big(\pred{P}_{ji}-P^\star_{ji}\big)^2\Big\rangle_{\mathcal{E},\,\mathcal{F}},
\qquad
L_{\mathrm{br},Q}\ \text{analogously}.
\label{eq:sfm-br}
\end{equation}

\paragraph{(7) Supervised thermal loading.}
(\code{loss.py:557--582}) With the tolerance-inflated rating
$S^{\max}_{\mathrm{eff}}=\sqrt{(S^{\max})^2+10^{-4}}$ and the worst-end convention
$\eta=\max(|S_{ij}|,|S_{ji}|)/S^{\max}_{\mathrm{eff}}$:
\begin{equation}
L_{\mathrm{therm}}=\Big\langle \logp\!\big((\pred\eta_e-\eta^\star_e)^2\big)\Big\rangle_{\mathcal{E},\,\mathcal{F}} .
\label{eq:sfm-loading}
\end{equation}

\paragraph{(8) Thermal-limit barrier with soft top-$k$ focus.}
(\code{loss.py:472--526}) This is the term that actually pushes \eqref{eq:therm}. Per rated
branch, with $\beta=8$ and $\varepsilon_{\mathrm{mean}}=0.05$:
\begin{equation}
\boxed{\;
v_e=\logp\!\big(\relu(\pred\eta_e-1)\big),
\qquad
L_{\mathrm{lim}}
=\Big\langle
\underbrace{\sum_{e\in k}\frac{e^{\beta v_e}}{\sum_{e'\in k}e^{\beta v_{e'}}}\,v_e}_{\text{soft-max over branches}}
+\;\varepsilon_{\mathrm{mean}}\,\langle v_e\rangle_k
\Big\rangle_{k} .
\;}
\label{eq:sfm-thermlim}
\end{equation}
The softmax weighting makes the term concentrate on the most overloaded branches of each
grid (a differentiable surrogate for $\max_e v_e$) while the $\varepsilon_{\mathrm{mean}}$
tail keeps a gradient on the rest.

\paragraph{(9) Physics-stress regulariser --- where angle limits enter.}
(\code{loss.py:741--758}, \code{stress\_features.py}) A 24-dimensional per-graph stress
vector is computed; its first $18$ dimensions are violation statistics
$(\text{mean},\max,\text{fraction}>10^{-3})$ of six families --- voltage, $\mathrm{KCL}_P$,
$\mathrm{KCL}_Q$, thermal, angle, and cycle/KVL imbalance --- with
\begin{align}
\text{voltage:}\quad & \relu(V^{\min}_i-\pred{V}_i)+\relu(\pred{V}_i-V^{\max}_i),
\label{eq:sfm-stress-v}\\
\text{thermal:}\quad & \relu\!\big(\max(|\pred{S}_{ij}|,|\pred{S}_{ji}|)/S^{\max}_{ij}-1\big),
\label{eq:sfm-stress-t}\\
\text{angle:}\quad & \relu\!\Big(\big|\pred\theta_i-\pred\theta_j-\phi_{ij}\big|\big/\Delta\theta^{\max}_{ij}-1\Big),
\label{eq:sfm-stress-a}\\
\text{KVL:}\quad & \Big|\textstyle\sum_{e\in c}\sigma_{ce}\,\big(\Delta\pred\theta_e-\phi_e\big)\Big|
\quad\text{over cycles } c ,
\label{eq:sfm-stress-kvl}
\end{align}
the remaining six dimensions being non-violation capacity-utilisation features. The
regulariser targets those 18 at zero:
\begin{equation}
\boxed{\;
L_{\mathrm{stress}}
=\Big\langle \logp\!\big(|\,\mathrm{stress}_{k,1:18}|\big)\Big\rangle_{k\in\mathcal{F}} .
\;}
\label{eq:sfm-stress}
\end{equation}
So the angle limits \eqref{eq:ang} and KVL are present in the objective, but only through
this pooled statistic --- three numbers per family per graph --- rather than as a per-branch
barrier of the form \eqref{eq:sfm-thermlim}. This is the precise sense in which GridSFM
handles angle limits \emph{indirectly}.

\paragraph{(10) Cost matching in log space.}
(\code{loss.py:636--646}, \code{730--735}) With the per-graph quadratic cost
\begin{equation}
\pred{C}_k=\!\!\sum_{g:\,b(g)=k}\!\!\Big(c_{2,g}\pred{P}_g^2+c_{1,g}\pred{P}_g+c_{0,g}\Big),
\qquad
C^\star_k \text{ the same formula at } P^\star_g \ \text{(detached)},
\label{eq:sfm-cost-def}
\end{equation}
the loss is a squared log-ratio, \emph{not} a minimisation of $\pred{C}$:
\begin{equation}
\boxed{\;
L_{\mathrm{cost}}
=\Big\langle\Big(\logp(\pred{C}_k)-\logp(C^\star_k)\Big)^{2}\Big\rangle_{k\in\mathcal{F}} .
\;}
\label{eq:sfm-cost}
\end{equation}
This is the important asymmetry with \eqref{eq:opf-obj}: the term is minimised at
$\pred{C}_k=C^\star_k$, so it penalises \emph{under}-cost as hard as over-cost. A free
minimisation of $\pred{C}$ would collapse to $\pred{P}_g\to P^{\min}_g$; matching the
optimal cost cannot. The $\log$ makes the term scale-free across grids whose costs differ
by orders of magnitude.

\paragraph{(11) Feasibility classification.}
(\code{loss.py:737--739}) On the \emph{whole} batch, including the synthetic infeasible
samples:
\begin{equation}
L_{\mathrm{feas}}=\operatorname{BCEWithLogits}\!\big(\,\pred{z}^{\mathrm{feas}},\ \text{feasible}\,\big)
=-\Big\langle y_k\log\sig(\pred z_k)+(1-y_k)\log\big(1-\sig(\pred z_k)\big)\Big\rangle_k .
\label{eq:sfm-feas}
\end{equation}
The head reads the stress vector \eqref{eq:sfm-stress-v}--\eqref{eq:sfm-stress-kvl} of the
model's \emph{own} prediction (\code{model.py:273--281}), so the classifier answers ``does
my predicted operating point look physically strained'' rather than inspecting the raw
instance alone.

\subsection{Total objective}

Assembling \code{loss.py:773--782} with its default multipliers:
\begin{equation}
\boxed{
\begin{aligned}
L_{\mathrm{GridSFM}}
={}& L_\theta+L_V+L_{P_g}+L_{Q_g}
\;+\;\lambda_{\mathrm{feas}}L_{\mathrm{feas}}
\;+\;\lambda_{\mathrm{cost}}L_{\mathrm{cost}}
\;+\;\lambda_{\mathrm{stress}}L_{\mathrm{stress}}\\
&+\lambda_{\mathrm{KCL},P}\,\logp\!\big(L_{\mathrm{KCL},P}\big)
+\lambda_{\mathrm{KCL},Q}\,\logp\!\big(L_{\mathrm{KCL},Q}\big)\\
&+\lambda_{\mathrm{br},P}\,\logp\!\big(L_{\mathrm{br},P}\big)
+\lambda_{\mathrm{br},Q}\,\logp\!\big(L_{\mathrm{br},Q}\big)
+\lambda_{\mathrm{therm}}L_{\mathrm{therm}}
+\lambda_{\mathrm{lim}}L_{\mathrm{lim}} ,
\end{aligned}}
\label{eq:sfm-total}
\end{equation}
\begin{gather*}
\lambda_{\mathrm{feas}}=0.1,\qquad
\lambda_{\mathrm{cost}}=1,\qquad
\lambda_{\mathrm{stress}}=0.1,\qquad
\lambda_{\mathrm{therm}}=1,\qquad
\lambda_{\mathrm{lim}}=5,\\
\lambda_{\mathrm{KCL},P}=\lambda_{\mathrm{KCL},Q}=\lambda_{\mathrm{br},P}=\lambda_{\mathrm{br},Q}=1 .
\end{gather*}
The outer $\logp$ on the four residual terms is a deliberate compressor: it bounds the
gradient contribution of an early-training KCL blow-up ($\partial\logp(L)/\partial L
=1/(1+L)$), so no single physics term can suppress the supervised signal. Note that the
logged diagnostics \code{parts['L\_kcl\_p']} etc.\ are the \emph{raw} pre-$\logp$ values,
while the contribution to \eqref{eq:sfm-total} is the compressed one.

% ======================================================================
\section{LUMINA}
\label{sec:lumina}

\subsection{The two artefacts, and which one answers which question}

Two LUMINA code drops are present, and they answer different questions:
\begin{itemize}
\item \code{LUMINA/src/lumina\_inference} --- an \textbf{inference-only} package: a 6-layer,
      128-channel, 2-head HGT with $0.244$ dropout over the node types $\{$\code{bus},
      \code{generator}, \code{load}, \code{shunt}$\}$ (\code{checkpoints/lumina\_config.json}),
      a dataset adapter, and a \code{Modeler} wrapper. No loss module, no optimiser, no
      training loop.
\item \code{lumina-sdk} --- the full SDK
      (\url{https://github.com/argonne-gridfm/lumina-sdk}, \code{main} at \code{f78cfae},
      2026-07-13), which \emph{does} ship the training stack: \code{lumina/model/opf/losses.py},
      \code{lumina/trainer/opf/trainer.py}, \code{lumina/evaluator/opf/evaluator.py}, and the
      configs under \code{configs/}.
\end{itemize}
Everything in this section is now line-traceable in the SDK. The result is more restrictive
than the paper-level description would suggest, so it is worth stating the headline first:
\emph{the released LUMINA training objective is a plain supervised regression loss. It
contains no power-balance term, no thermal term, no cost term, and no Lagrangian.}

\subsection{Predictions and architectural (hard) box constraints}

Every heterogeneous architecture in the SDK (\code{OPFHeteroGNN}, \code{RGAT}, \code{HEAT},
\code{HGT}) emits two channels per bus and two per generator,
\begin{equation}
\pred{y}_{\mathrm{bus}}=(\pred\theta,\pred{V}_m)\in\Real^{N_b\times 2},
\qquad
\pred{y}_{\mathrm{gen}}=(\pred{P}_g,\pred{Q}_g)\in\Real^{N_g\times 2},
\label{eq:lum-heads}
\end{equation}
and each caches the raw limit columns of the \emph{unnormalised} inputs before the first
projection, then rescales through a sigmoid when \code{minmax\_scaling=True}
(\code{lumina/model/opf/hetero\_model.py:186--196, 333--343, 546--554, 712--722}):
\begin{align}
\pred{V}_i &= V^{\min}_i+\big(V^{\max}_i-V^{\min}_i\big)\,\sig\!\big(z_{V,i}\big),
&& V^{\min}=x^{\mathrm{bus}}_{:,1},\; V^{\max}=x^{\mathrm{bus}}_{:,2},
\label{eq:lum-v}\\
\pred{P}_g &= P^{\min}_g+\big(P^{\max}_g-P^{\min}_g\big)\,\sig\!\big(z_{P,g}\big),
&& P^{\min}=x^{\mathrm{gen}}_{:,2},\; P^{\max}=x^{\mathrm{gen}}_{:,3},
\label{eq:lum-p}\\
\pred{Q}_g &= Q^{\min}_g+\big(Q^{\max}_g-Q^{\min}_g\big)\,\sig\!\big(z_{Q,g}\big),
&& Q^{\min}=x^{\mathrm{gen}}_{:,5},\; Q^{\max}=x^{\mathrm{gen}}_{:,6},
\label{eq:lum-q}
\end{align}
the column indices matching the loader layout
$x^{\mathrm{bus}}=[\,\mathrm{base\_kv},V^{\min},V^{\max},\mathrm{onehot}(\mathrm{bus\_type})\,]\in\Real^{7}$
and
$x^{\mathrm{gen}}=[\,\mathrm{mbase},P_g,P^{\min},P^{\max},Q_g,Q^{\min},Q^{\max},V_g,c_2,c_1,c_0\,]\in\Real^{11}$
(\code{docs/opfdata\_schema.md}, \code{dataset/opf/opf\_dataset.py:700--710}). Since
$0<\sig(z)<1$ strictly, \eqref{eq:vlim} and \eqref{eq:glim} hold with strict inequality for
every prediction --- a genuine architectural guarantee, the same mechanism as GridFM's
\eqref{eq:gridfm-vbound}--\eqref{eq:gridfm-pbound}, extended to $Q_g$. The trainer defaults
to \code{minmax\_scaling=True} (\code{trainer/opf/trainer.py:141}), so this is the normal
training and inference regime. Three qualifications:
\begin{itemize}
\item The bus \emph{angle} channel is passed through untransformed, so nothing bounds
      $\pred\theta$ or the angle differences \eqref{eq:ang}.
\item With \code{minmax\_scaling=False} no constraint is enforced anywhere; the guarantee is
      a property of the call, not of the checkpoint.
\item The model forward signatures accept \code{edge\_attr\_dict}, but the trainer never
      supplies it (\code{trainer.py:1387} calls the model with \code{x\_dict} and
      \code{edge\_index\_dict} only), and \code{HGT} --- the released checkpoint's
      architecture --- ignores it regardless, since \code{HGTConv} is invoked as
      \code{conv(x\_dict, edge\_index\_dict)} (\code{hetero\_model.py:844}). Branch
      $r,x,b,\tau,\phi$ and the ratings $S^{\max}$ therefore never enter the hetero training
      path at all. (\code{HEAT} and the \code{gat} backend do project edge attributes ---
      \code{hetero\_model.py:695--700} --- but only if a caller passes them.)
\end{itemize}

\subsection{The actual training loss: per-node-type supervised regression}

\code{ACOPFLossFunction} (\code{lumina/model/opf/losses.py:83--348}) is the whole objective.
Targets are the solver solution channels, $y_{\mathrm{bus}}=(\theta^\star,V^\star_m)$ and
$y_{\mathrm{gen}}=(P^\star_g,Q^\star_g)$, extracted as \code{batch[node\_type].y}. Rows with
any non-finite target are dropped by a finite mask $\mathcal{V}_t$
(\code{losses.py:184--193}), and the per-type losses are combined with node weights
$w_t$ (default $w_{\mathrm{bus}}=w_{\mathrm{gen}}=1$):
\begin{equation}
\boxed{\;
L_{\mathrm{LUMINA}}
=\sum_{t\in\{\mathrm{bus},\,\mathrm{gen}\}} w_t\;
\ell\!\left(\pred{y}_t[\mathcal{V}_t],\;y_t[\mathcal{V}_t]\right),
\;}
\label{eq:lum-total}
\end{equation}
where $\ell$ is \emph{one} of five element-wise choices, selected by \code{loss\_type}
(\code{configs/loss/mse.yaml}, \code{docs/tutorials/training.md}), reduced by mean over all
$N$ surviving elements $\times$ channels:
\begin{align}
\text{\code{mse}:}\quad & \ell=\tfrac{1}{N}\textstyle\sum_n (\pred{y}_n-y_n)^2,
\label{eq:lum-mse}\\
\text{\code{mae}:}\quad & \ell=\tfrac{1}{N}\textstyle\sum_n |\pred{y}_n-y_n|,
\label{eq:lum-mae}\\
\text{\code{rmse}:}\quad & \ell=\tfrac{1}{N}\textstyle\sum_n \sqrt{(\pred{y}_n-y_n)^2+\epsilon},
\qquad \epsilon=10^{-8},
\label{eq:lum-rmse}\\
\text{\code{mape}:}\quad & \ell=\tfrac{1}{N}\textstyle\sum_n \dfrac{|\pred{y}_n-y_n|}{|y_n|+\epsilon},
\label{eq:lum-mape}\\
\text{\code{smooth\_l1}:}\quad & \ell=\tfrac{1}{N}\textstyle\sum_n
\begin{cases}
\dfrac{(\pred{y}_n-y_n)^2}{2\beta}, & |\pred{y}_n-y_n|<\beta,\\[4pt]
|\pred{y}_n-y_n|-\dfrac{\beta}{2}, & \text{otherwise,}
\end{cases}
\qquad \beta=1 .
\label{eq:lum-huber}
\end{align}
With the default \code{loss\_type: mse}, \eqref{eq:lum-total} is exactly the generic
supervised term \eqref{eq:sup} restricted to the four predicted channels --- nothing more.
Note that \eqref{eq:lum-rmse}, despite its name, takes the square root \emph{element-wise
before} averaging, so it is an $\epsilon$-smoothed MAE, not
$\sqrt{\langle(\pred{y}-y)^2\rangle}$; the two differ by Jensen's inequality and have
different gradients.

\paragraph{No cost, no KCL, no thermal term.} A \code{grep} for \code{cost} over
\code{losses.py} and \code{trainer.py} returns nothing: the cost coefficients
$c_2,c_1,c_0$ occupy generator input columns 8--10 and are read only by the evaluator
helper \code{extract\_generation\_costs\_from\_batch}. The solver's optimal objective is
stored in the dataset (\code{objective} field, \code{docs/opfdata\_schema.md:263}) and is
not used in training. So cost reaches the parameters exactly as in GridFM --- only through
the labels $P^\star_g$ --- and \eqref{eq:kcl-p}--\eqref{eq:kcl-q}, \eqref{eq:therm} and
\eqref{eq:ang} have no gradient path whatsoever.

\subsection{The physics-informed loss is a skeleton}

\code{PhysicsInformedLoss} (\code{losses.py:267--348}) subclasses the above and adds
\begin{equation}
L_{\mathrm{PI}} = L_{\mathrm{LUMINA}} + w_{\mathrm{phys}}\,\mathcal{P},
\qquad
\mathcal{P}=\sum_{c\,\in\,\mathcal{C}} \pi\big(v_c(\pred{x})\big),
\qquad
\mathcal{C}\subseteq\{\text{power\_flow},\text{line\_limits}\},
\label{eq:lum-pi}
\end{equation}
with three selectable penalty forms (\code{losses.py:324--334})
\begin{equation}
\pi_{\mathrm{quad}}(v)=\big\langle v^2\big\rangle,
\qquad
\pi_{\mathrm{abs}}(v)=\big\langle |v|\big\rangle,
\qquad
\pi_{\mathrm{barrier}}(v)=-\big\langle \log\max(-v,\epsilon)\big\rangle .
\label{eq:lum-penalty}
\end{equation}
But $v_c$ is supplied by an injected \code{constraint\_computer} that defaults to
\code{None}, in which case \code{\_compute\_physics\_penalty} returns the constant $0$
(\code{losses.py:314--316}); \code{set\_constraint\_computer} is never called anywhere in
the SDK, the examples, or the tests, and the trainer instantiates \code{OPFLossManager}
(hence \code{ACOPFLossFunction}), never \code{PhysicsInformedLoss}
(\code{trainer.py:1689, 2335}). Both the README and the changelog call it a ``skeleton''.
Therefore $\mathcal{P}\equiv 0$ and \eqref{eq:lum-pi} collapses to \eqref{eq:lum-total} in
every shipped code path.

\subsection{The augmented Lagrangian was removed before release}

The generic augmented-Lagrangian objective, with $h$ the equality residual (power balance)
and $g$ the inequality residual, is
\begin{equation}
L_{\mathrm{AL}}
= L_{\mathrm{sup}}(\pred{x},x^\star)
\;+\;\lambda^\top h(\pred{x})
\;+\;\frac{\rho}{2}\big\|h(\pred{x})\big\|^2
\;+\;\mu^\top \relu\!\big(g(\pred{x})\big)
\;+\;\frac{\rho}{2}\big\|\relu\!\big(g(\pred{x})\big)\big\|^{2} ,
\label{eq:lum-al}
\end{equation}
with $(\lambda,\mu,\rho)$ updated on a slower schedule than $\phi$, e.g.
$\lambda\leftarrow\lambda+\rho\,h(\pred{x})$. Its distinction from \eqref{eq:kkt} is
structural: in IPOPT, $\lambda$ and $\mu$ are unknowns solved \emph{jointly with} $x$ for
\emph{each instance} and certify optimality at convergence; in \eqref{eq:lum-al} they are
training-time hyper-state shared \emph{across} instances, shaping $\phi$ so that predictions
tend to have small violations, and nothing in \eqref{eq:lum-al} is evaluated at inference.

\eqref{eq:lum-al} is included here for contrast only. It is \textbf{not} the released
objective: the SDK changelog records, under \code{[Unreleased]} $\to$ \emph{Changed},
``\emph{Removed unreleased Lagrangian settings from Frontier-facing configs and removed
Lagrangian loss options from the Frontier training CLI}'' (\code{CHANGELOG.md:17}), and a
\code{grep -i lagrang} over the entire SDK matches that changelog line and nothing else ---
no loss class, no config key, no CLI flag. The Lagrangian variants studied in the papers
are thus not reproducible from the public SDK, and the claim that LUMINA's thermal and
angle inequalities are ``handled in the augmented-Lagrangian term'' cannot be supported by
any released artefact.

\subsection{Constraints appear only as evaluation metrics --- and only the box ones}

\code{ACOPFConstraintEvaluator} (\code{evaluator/opf/evaluator.py:194--373}) computes, with
no gradient path,
\begin{align}
\mathcal{V}_V &= \Big\langle \relu\!\big(V^{\min}_i-\pred{V}_i\big)+\relu\!\big(\pred{V}_i-V^{\max}_i\big)\Big\rangle_i,
\label{eq:lum-viol-v}\\
\mathcal{V}_{P} &= \Big\langle \relu\!\big(P^{\min}_g-\pred{P}_g\big)+\relu\!\big(\pred{P}_g-P^{\max}_g\big)\Big\rangle_g,
\qquad
\mathcal{V}_{Q}\ \text{analogously},
\label{eq:lum-viol-pq}\\
\mathcal{V}_{\mathrm{total}} &= \mathcal{V}_V+\mathcal{V}_P+\mathcal{V}_Q .
\label{eq:lum-viol-tot}
\end{align}
\code{evaluate\_all\_constraints} delegates \emph{only} to these bound checks; its
docstring states that ``power-flow and line-flow evaluations are planned but not yet
active'', and the aggregation of \code{real\_power\_flow\_violations},
\code{reactive\_power\_flow\_violations} and \code{line\_flow\_violations} into a total is
commented out (\code{evaluator.py:404--408}). So the KCL and thermal violations that
\eqref{eq:sfm-kcl} and \eqref{eq:sfm-thermlim} penalise in GridSFM are, in LUMINA, neither
trained nor measured.

There is a further subtlety worth noting: under the default
\code{minmax\_scaling=True}, \eqref{eq:lum-v}--\eqref{eq:lum-q} make
\eqref{eq:lum-viol-v}--\eqref{eq:lum-viol-pq} identically zero by construction. The only
violation metric the SDK reports is therefore the one the architecture already guarantees,
and a reported $\mathcal{V}_{\mathrm{total}}=0$ is evidence about the output layer, not
about feasibility of the operating point.

\paragraph{Model selection does not see constraints either.} The checkpoint monitor is
\code{val/score} (\code{configs/config.yaml:44}), and
\begin{equation}
\texttt{val/score} \;=\; \texttt{val/loss/objective} \;=\; L_{\mathrm{LUMINA}}
\label{eq:lum-score}
\end{equation}
exactly (\code{trainer.py:1334--1340}). The config keys \code{score\_alpha},
\code{log\_normalized\_violation}, \code{violation\_eval\_p} are parsed and stored, and the
sub-sampled ``violation evaluation'' pass (\code{trainer.py:1854--1960}) in fact recomputes
only the regression loss; no violation term is folded into the score. Best-model selection
is therefore purely regression-driven.

\subsection{Implementation notes observed while tracing the SDK}

These are incidental to the formulas but affect whether the code above runs as written, on
\code{main} at \code{f78cfae}:
\begin{enumerate}
\item \code{trainer.py:1692} (and \code{2338}) construct
      \code{OPFLossManager(loss\_type=..., device=..., log\_normalized\_violation=...)}.
      \code{OPFLossManager} forwards \code{**kwargs} to \code{ACOPFLossFunction}
      (\code{losses.py:377}), whose \code{\_\_init\_\_} accepts only
      \code{loss\_type, node\_weights, reduction, epsilon, beta} and no \code{**kwargs}
      (\code{losses.py:108--115}) --- so this raises
      \code{TypeError: \_\_init\_\_() got an unexpected keyword argument
      'log\_normalized\_violation'} at trainer construction. Verified statically from the
      AST; not executed here (no \code{torch} in this environment).
\item \code{create\_combined\_loss} and \code{create\_robust\_loss}
      (\code{losses.py:535--562}) pass \code{loss\_type} as a \emph{dict}
      (e.g.\ \code{\{'mse': 0.7, 'smooth\_l1': 0.3\}}), which fails the
      \code{loss\_type not in valid\_types} check at \code{losses.py:125}. The advertised
      weighted-combination loss is therefore unreachable; \eqref{eq:lum-total} admits one
      $\ell$ at a time.
\item \code{ACOPFLossFunction.compute\_loss} (\code{losses.py:219--230}) computes
      \code{masked\_predictions}/\code{masked\_targets} and then calls
      \code{self.forward(predictions, targets)} with the \emph{unmasked} pair. The
      \code{y\_mask} path is dead in that method (\code{OPFLossManager.compute\_loss}
      does use the masked pair).
\end{enumerate}
None of these change \eqref{eq:lum-total}; items 1 and 2 mean the published
\code{main} needs a patch before the documented training commands run.

% ======================================================================
\section{Side-by-side summary}

\renewcommand{\arraystretch}{1.25}
\begin{longtable}{@{}p{0.20\textwidth}p{0.25\textwidth}p{0.25\textwidth}p{0.22\textwidth}@{}}
\toprule
\textbf{OPF element} & \textbf{GridFM} & \textbf{GridSFM} & \textbf{LUMINA} (SDK)\\
\midrule
\endhead
Predicts & $V_m,V_a,P_g$; $Q_g$ \emph{derived} from Q-balance \eqref{eq:gridfm-qg}
         & $\theta,V,P_g,Q_g$, flows, feasibility logit
         & $V_a,V_m,P_g,Q_g$ \\
KCL equality
         & layered residual penalty \eqref{eq:gridfm-layered}; \code{PBE} \eqref{eq:gridfm-pbe}; exact in $Q$ at PV/REF
         & capacity-normalised worst-bus-tilted penalty \eqref{eq:sfm-kcl}
         & \textbf{no} term; not even evaluated \\
$V$ bounds
         & \textbf{hard} sigmoid \eqref{eq:gridfm-vbound} (OPF only)
         & \textbf{hard} leaky clip \eqref{eq:sfm-v} ($\alpha{=}0$)
         & \textbf{hard} sigmoid \eqref{eq:lum-v} \\
$P_g$ bounds
         & \textbf{hard} sigmoid \eqref{eq:gridfm-pbound} (OPF only)
         & \textbf{hard} leaky clip \eqref{eq:sfm-pg}
         & \textbf{hard} sigmoid \eqref{eq:lum-p} \\
$Q_g$ bounds
         & soft hinge over violators \eqref{eq:gridfm-qviol}, weight $10^{-3}$
         & \textbf{hard} leaky clip \eqref{eq:sfm-qg}
         & \textbf{hard} sigmoid \eqref{eq:lum-q} \\
Thermal limit
         & evaluation only (\code{test\_step})
         & soft top-$k$ barrier \eqref{eq:sfm-thermlim} ($\lambda{=}5$) + supervised loading \eqref{eq:sfm-loading}
         & \textbf{no} term; ``planned, not yet active'' in evaluator \\
Angle limit
         & evaluation only
         & pooled stress statistic \eqref{eq:sfm-stress-a} inside \eqref{eq:sfm-stress}
         & \textbf{no} term; $\theta$ head unscaled \\
Generation cost
         & \textbf{no} loss term; optimality gap \eqref{eq:gridfm-gap} at test time
         & explicit log-cost \emph{matching} \eqref{eq:sfm-cost} + dispatch imitation
         & \textbf{no} term; $c_0,c_1,c_2$ inputs unused in training \\
Feasibility
         & --- & BCE head \eqref{eq:sfm-feas} + synthetic infeasible samples & --- \\
KKT solve at inference & no & no & no \\
Objective & \eqref{eq:gridfm-instance} & \eqref{eq:sfm-total} & \eqref{eq:lum-total}: one of MSE / MAE / RMSE / MAPE / Huber \\
Warm start to exact solver & possible & demonstrated (PowerModels/IPOPT) & possible in principle \\
\bottomrule
\end{longtable}

% ======================================================================
\section{Reading of the whole picture}

Classical OPF solves one optimisation problem per sample. These three solve many offline,
then fit \eqref{eq:map} to the solutions, and answer in one forward pass. Their training
signals decompose into six pedagogically distinct instructions, and the strength of each is
what the formulas above make precise:
\begin{center}
\renewcommand{\arraystretch}{1.2}
\begin{tabular}{@{}p{0.30\textwidth}p{0.20\textwidth}p{0.42\textwidth}@{}}
\toprule
Mechanism & Strength & Formula \\
\midrule
``be close to the solver solution'' & soft, but the anchor of everything
  & \eqref{eq:gridfm-masked}, \eqref{eq:sfm-Ltheta}--\eqref{eq:sfm-Lqg}, \eqref{eq:lum-total} \\
``reduce the power-balance mismatch'' & soft
  & \eqref{eq:gridfm-layered}, \eqref{eq:gridfm-pbe}, \eqref{eq:sfm-kcl} \\
``never leave this interval'' & \textbf{hard} (architectural)
  & \eqref{eq:gridfm-vbound}--\eqref{eq:gridfm-pbound}, \eqref{eq:sfm-pg}--\eqref{eq:sfm-v}, \eqref{eq:lum-v}--\eqref{eq:lum-q} \\
``do not overload a line'' & soft
  & \eqref{eq:sfm-thermlim} \\
``match the optimal cost'' & soft, two-sided
  & \eqref{eq:sfm-cost} \\
``recognise an infeasible instance'' & classification
  & \eqref{eq:sfm-feas} \\
\bottomrule
\end{tabular}
\end{center}

Three conclusions follow directly from the formulas rather than from intuition.

\textbf{First}, only constraints realised as output transformations are guaranteed:
\eqref{eq:vlim} and \eqref{eq:glim} in GridSFM and in LUMINA-with-scaling, and
\eqref{eq:vlim} plus the $P_g$ box in GridFM's OPF task. Everything else --- KCL, thermal,
angle --- is a penalty and can be violated at inference by an amount no term in
\eqref{eq:sfm-total} bounds a priori.

\textbf{Second}, ``minimising generation cost'' means something different in each model.
GridFM has no cost gradient at all: cost enters exclusively through the labels
\eqref{eq:gridfm-masked}, and \eqref{eq:gridfm-gap} is a report card. GridSFM has a cost
gradient, but it is a \emph{matching} gradient \eqref{eq:sfm-cost} that is equally unhappy
about a cheaper-than-optimal dispatch --- which is the correct design, since an
unconstrained $\min\pred{C}$ under a soft KCL penalty has the degenerate solution
$\pred{P}_g\to P^{\min}_g$. Neither model performs a descent on \eqref{eq:opf-obj} subject
to \eqref{eq:kcl-p}--\eqref{eq:ang} at inference.

\textbf{Third}, the three models occupy three distinct rungs of constraint awareness, and
the gap is larger than the papers suggest. GridSFM penalises every family
\eqref{eq:sfm-total} and matches the cost \eqref{eq:sfm-cost}. GridFM penalises power
balance \eqref{eq:gridfm-layered} and the reactive limits \eqref{eq:gridfm-qviol}, and
evaluates the rest. The released LUMINA objective \eqref{eq:lum-total} is a single
regression loss over four channels: its only constraint mechanism is the sigmoid box
\eqref{eq:lum-v}--\eqref{eq:lum-q}, the physics penalty \eqref{eq:lum-pi} is a skeleton
whose penalty term is identically zero, the Lagrangian options were removed before release
(\code{CHANGELOG.md:17}), and even the evaluator measures only the box constraints the
output layer already guarantees. So when comparing feasibility metrics across the three,
note that they are not three implementations of the same idea: GridSFM is trained against
the violations, GridFM against some of them, and LUMINA against none. Accordingly,
\eqref{eq:lum-al} should be cited as a contrast case or as paper-reported behaviour, never
as a property of the released checkpoint or SDK.

\end{document}
